\documentclass[11pt,a4paper]{article}
\usepackage[margin=22mm]{geometry}
\usepackage{amsmath,amssymb,amsthm}
\newtheorem{proposition}{Proposition}
\title{A common action for a group and its automorphisms}
\author{Human Compression project}\date{}
\begin{document}\maketitle
The acquired regular-polytope enumerator requires permutation groups
$G\mathrel{\triangleleft}A$ in the same action, with
$|A:C_A(G)|=|\operatorname{Aut}(G)|$. A separately named permutation
model of $\operatorname{Aut}(G)$ does not by itself supply this input.
We use the following elementary construction to connect already
identified models to the acquired enumerator.

\begin{proposition}
Let $H$ be a finite permutation group and $a=|\operatorname{Aut}(H)|$.
Either of the following constructions supplies a faithful copy $G$ of
$H$ and a conjugating group $A$ inducing every automorphism of $G$.
\begin{enumerate}
\item Begin with $A=H$ in its given action. Adjoin permutations which
normalize $H$. Accept this action when $|A:C_A(H)|=a$ and put $G=H$.
\item Let $I$ be the complete set of involutions of $H$. Let $G$ be the
image of the conjugation action of $H$ on $I$, and let $A$ be the image
of the action of $\operatorname{Aut}(H)$ on $I$. Accept these actions
when both are faithful. Then $G\mathrel{\triangleleft}A$ and
$|A:C_A(G)|=a$.
\end{enumerate}
\end{proposition}
\begin{proof}
In the first construction, the supplied permutations normalize $H$, so
the resulting group $A$ does too. Conjugation gives a homomorphism
$A\to\operatorname{Aut}(H)$ with kernel $C_A(H)$. Equality of the
finite orders implies that its image is all of $\operatorname{Aut}(H)$.

For the second construction, automorphisms preserve element order, so
$I$ is invariant under both actions. Write
$\iota:H\to\operatorname{Sym}(I)$ and
$\alpha:\operatorname{Aut}(H)\to\operatorname{Sym}(I)$ for their
homomorphisms. Their images are $G$ and $A$. Faithfulness of $\iota$
identifies $H$ with $G$; this is a checked condition, not an assumption
based on the group name. For $h\in H$ and $\phi\in\operatorname{Aut}(H)$,
the conjugation actions satisfy
\[
 \alpha(\phi)^{-1}\iota(h)\alpha(\phi)=\iota(\phi(h)),
\]
with permutations acting on the right. Hence $G$ is normal in $A$,
and the induced automorphisms of $G$ are precisely the transports of
all automorphisms of $H$. If $\alpha(\phi)$ centralizes $G$, the
displayed equality and injectivity of $\iota$ give $\phi(h)=h$ for
every $h$. Thus $C_A(G)$ is trivial and $|A:C_A(G)|=a$.
\end{proof}

In the implementation, a proposed normalizing permutation is tested
on every generator of $H$. Inclusion $H^c\leq H$ implies equality
because conjugation preserves the finite order. Candidate overgroups
are only a way of obtaining a smaller faithful action; their names and
orders are never evidence for the automorphism condition. If none
suffices, the implementation forms $I$ from all involution conjugacy
classes and checks its cardinality against their sizes. The orders of
both action images certify faithfulness. It additionally checks
normality and the centralizer quotient directly.

The generator images of $H$ under $\iota$ retain the connection to the
original named model. The acquired enumerator then applies unchanged
to $(G,A)$, and an isomorphism preserves generation, string relations,
all subset intersections and equivalence under automorphisms. Its
ordered classifications therefore transfer to $H$. Completeness still
requires a complete returned enumeration; successful action setup alone
is not a polytope census.

The use of involutions as a common set is borrowed from the already
reviewed exceptional Coxeter computations in F006. The enumeration,
intersection criterion, rank bound and exact GAP operations are
acquired from the reviewed L071 application. This component provides
the compatibility argument and explicit checked interface, not a new
classification theorem or an open-conjecture solution.
\end{document}
